\section{From a small audit to a ranking}
\label{sec:setup}

This section turns the question ``which orderings survive a human reader?'' into something we can
compute. We define a model's accuracy under the reader verdict, show why 400 labels cannot settle
a ranking by themselves, state the one assumption that lets models share labels, and explain when
an ordering counts as settled and with what confidence. Proofs are in Appendix~\ref{app:proofs}.

\subsection{What an ordering claims}

Each model's responses are spread over the eight cells. Let $s_{mc}$ be the \emph{share} of model
$m$'s responses that fall in cell $c$; shares are known exactly, because every verifier reads every
response. Let $p_{mc}$ be the fraction of those responses with reader verdict $Y^* = 1$. These
fractions are what we do not know. Model $m$'s accuracy under the reader verdict, and the gap
between models $j$ and $k$, are
\begin{equation}
A_m \;=\; \sum_{c} s_{mc}\, p_{mc}, \qquad \Delta_{jk} = A_j - A_k .
\label{eq:accuracy}
\end{equation}
An \emph{ordering} $j \succ k$ is the claim that $\Delta_{jk} > 0$: model $j$ would score higher
than model $k$ if a reader had marked every response. The \NumModels{} models form
\NumOrderings{} pairs, and for each pair we ask whether either ordering can be established.

The difficulty is coverage. The 400 labels leave most combinations of model and cell with no label
at all, and for those the only safe statement is $p_{mc} \in [0,1]$. Each such cell widens the range
of possible gaps by its share, and when the widening exceeds the gap itself the ordering stays open.
The discordant cells are small, so leaving them open costs little. The concordant cells hold almost
every response, so they must be pinned down tightly, and that requires letting models share labels.

\subsection{Letting models share labels}

Sharing labels between models, which statisticians call \emph{pooling}, is an assumption, not a
measurement, and we state it with a single number.

\begin{assumption}[Bounded heterogeneity, $H(\varepsilon)$]
\label{ass:h}
For each concordant cell $c$, the values of $p_{gc}$ across the audit sub-populations $g$ that
make up that cell lie in an interval of width at most $\varepsilon$.
\end{assumption}

In words: where all three verifiers agree, the rate at which a careful reader would overrule them
may differ from model to model, but by no more than $\varepsilon$. A sub-population is usually a
model; where the audit split a model's responses into groups (by subject and by the boxed
verifier's verdict), each group counts separately, which makes the assumption slightly stronger.
The assumption says nothing about the overrule rate itself, which the audit estimates, only about
how much it varies.

Setting $\varepsilon = 0$ is stronger than it looks. Each $p_{gc}$ is a fraction of a finite set of
responses, so two models with the same underlying chance of being overruled will still, by luck,
end up with slightly different fractions. $H(0)$ is therefore a condition on the actual responses,
not a consequence of equal rates. We report every result as a function of $\varepsilon$ and mark
$\varepsilon = \EpsPrereg$, fixed in our analysis plan as a reference. That value is a judgement,
not a measurement, and it was set after the audit's labels existed (Appendix
Table~\ref{tab:chronology}), which is why no count in this paper appears without the curve around
it.

\subsection{Certified orderings and breakdown values}

The idea is simple. In each concordant cell the audit gives a confidence interval for the shared
rate, and under the assumption every model's own rate lies within $\varepsilon$ of it. In each
discordant cell a model's rate may be anything. Taking the least and the most favourable values in
every cell gives the smallest and the largest gap two models could have. Formally, let $\tilde p_c$
be the pooled fraction of responses with $Y^* = 1$ in concordant cell $c$,
weighted as the audit sampled the sub-populations, and let $[L_c, U_c]$ be a confidence interval
for it, built in Section~\ref{sec:inference} so that it does not assume the models are alike. For a
discordant cell, let $[l_{mc}, u_{mc}]$ be an interval for model $m$ alone, which is $[0,1]$ when
that model has no labels in the cell.

\begin{proposition}[Identified range]
\label{prop:range}
Fix a pair $(j,k)$ and a concordant cell $c$. Under Assumption~\ref{ass:h} and on the event
$\tilde p_{c} \in [L_c, U_c]$, the pair $(p_{jc}, p_{kc})$ lies in
\[
R_c(\varepsilon) \;=\; \bigl\{ (x,y) \in [\ell_c, h_c]^2 \;:\; |x - y| \le \varepsilon \bigr\},
\qquad \ell_c = \max(0, L_c - \varepsilon), \quad h_c = \min(1, U_c + \varepsilon).
\]
Consequently $\Delta_{jk}$ lies in the interval obtained by minimising and maximising
Equation~\eqref{eq:accuracy} separately in each cell, over $R_c(\varepsilon)$ for concordant
cells and over the per-model intervals for discordant ones.
\end{proposition}

We call this interval the \emph{identified range} of $\Delta_{jk}$: every gap between the two
models that is consistent with the verdicts, the labels and the assumption. Appendix~\ref{app:proofs}
proves the proposition and solves the optimisation in closed form, and Appendix~\ref{app:tests}
checks both against brute force. At $\varepsilon = 0$ the range is exact; for $\varepsilon > 0$ it
is slightly wider than necessary, which can only make conclusions more cautious.

An ordering $j \succ k$ is \emph{certified} at $\varepsilon$ when the whole identified range of
$\Delta_{jk}$, widened by any sampling term from Section~\ref{sec:inference}, lies above zero. A
larger $\varepsilon$ gives a wider range, so an ordering certified at one tolerance is certified at
every smaller one, which makes the following well defined.

\begin{definition}[Breakdown value]
\label{def:breakdown}
$\varepsilon^*_{jk} = \sup\{\varepsilon \in [0, 1] : j \succ k \text{ is certified at }
\varepsilon\}$, with $\varepsilon^*_{jk} = 0$ when the ordering is not certified at
$\varepsilon = 0$.
\end{definition}

The breakdown value is how much model-to-model variation an ordering can absorb before it stops
being certified. Reporting it for every ordering turns the dependence on Assumption~\ref{ass:h} into
a number that a reader can hold against their own view. Certification is a strict inequality, so
the supremum need not be attained, and because the range is cautious for $\varepsilon > 0$, each
reported value understates the true one.

\subsection{Why three verifiers rather than one}

\begin{proposition}[A common error mechanism cannot reorder]
\label{prop:oneverifier}
Suppose one verifier is used, and suppose its predictive values are common to all models, so
that $\Pr(Y^* = 1 \mid v = 1) = \pi_1$ and $\Pr(Y^* = 1 \mid v = 0) = \pi_0$ regardless of $m$.
Then $\Delta_{jk} = (\pi_1 - \pi_0)(a_j - a_k)$, where $a_m$ is the rate at which the verifier
credits model $m$.
\end{proposition}

Every gap is scaled by the same factor, so a verifier whose mistakes are alike for all models
cannot swap one pair while leaving another in place. The swaps that Section~\ref{sec:generality}
reports therefore need mistakes that differ between models, and correcting a single verifier with
error rates shared by all models would assume away exactly what we study. Sharing labels only where
verifiers with different extraction rules agree leaves the disagreement cells free to behave
differently for each model.

\subsection{Confidence}
\label{sec:inference}

\paragraph{An interval that does not assume the models are alike.}
The pooled count in a concordant cell adds up labels from models whose true rates may differ, and
the labels were drawn without replacement within each stratum of the audit. A textbook binomial
interval assumes one common rate and sampling with replacement; using it here would quietly build
the sameness that Assumption~\ref{ass:h} makes explicit into the confidence statement. We use a
bound that holds without either condition.

\begin{lemma}[Design-valid pooled bound]
\label{lem:kl}
Let $x_c$ be the total number of audited responses with $Y^* = 1$ in cell $c$, summed over
sub-populations, with $n_c$ labels in total and sampling weights $w_{gc} = n_{gc}/n_c$. Let
$\tilde p_c = \sum_g w_{gc}\, p_{gc}$. Then, conditional on the realised number of labels in each
sub-population, for $0 < t \le 1 - \tilde p_c$ and $0 < t \le \tilde p_c$ respectively (larger $t$
make the events empty), with $\mathrm{KL}$ as defined in Appendix~\ref{app:proofs},
\[
\Pr\!\left(\frac{x_c}{n_c} \ge \tilde p_c + t\right) \le e^{-n_c \mathrm{KL}(\tilde p_c + t \,\|\, \tilde p_c)},
\qquad
\Pr\!\left(\frac{x_c}{n_c} \le \tilde p_c - t\right) \le e^{-n_c \mathrm{KL}(\tilde p_c - t \,\|\, \tilde p_c)} .
\]
\end{lemma}

The proof uses two standard facts about moment generating functions: sampling without
replacement is dominated by sampling with replacement, and a sum of independent Bernoulli
variables with different means is dominated by a single Bernoulli at their weighted mean.
Inverting the bound gives $[L_c, U_c]$ with at least nominal coverage in finite samples.

\paragraph{Two targets: these problems, and questions like them.}
The \emph{primary target} is $\Delta_{jk}$ on the \NumItems{} problems the leaderboard actually
used, because that is what a leaderboard reports: a score on a fixed set of questions. For this
target the only uncertainty comes from the audit, the only assumption is Assumption~\ref{ass:h}, and
the guarantee holds at the audit's actual size, with no large-sample approximation. Every headline number in the paper is of this kind.
A reader may also want a claim about questions like these in general. MATH-Hard is a curated
benchmark, not a random draw from a defined population, so such a claim has to treat the problems
as if they were one, which adds an uncertainty that no amount of labelling can remove. We bound it
with a margin for item sampling (a paired half-width) built from quantities the identified range
already supplies; it relies on a central limit approximation, and a distribution-free alternative is reported beside it
(Appendix~\ref{app:inference}). This \emph{secondary target} is analysed in full, and its numbers are
never mixed with the primary ones. The order of the two departs from our analysis plan, which first
named the secondary target as its headline; we made the change after the audit's results for both
targets were known and recorded it as an amendment to the plan (Appendix Table~\ref{tab:chronology}).

\paragraph{All orderings at once.}
Someone reading a leaderboard reads all \NumOrderings{} orderings together, so every headline count
is simultaneous: the audit's error level is divided over every parameter on the roster
(Appendix~\ref{app:inference}), and for the secondary target the sampling term also receives a
Bonferroni correction, which divides its error level among the orderings. Rank intervals follow directly from the certified set
(Appendix~\ref{app:proofs}).

\begin{proposition}[Coverage]
\label{prop:coverage}
Under Assumption~\ref{ass:h}, all claims about the primary target hold simultaneously with
probability at least $1 - \alpha_{\mathrm{audit}}$ at any sample size, with no large-sample
approximation. Claims about the
secondary target hold with probability at least $1 - \alpha_{\mathrm{audit}} -
\alpha_{\mathrm{sample}}$, asymptotically in the number of items.
\end{proposition}

The event behind the primary guarantee is that every concordant interval covers its target, and
that event does not involve $\varepsilon$. On it, the orderings certified at every tolerance for
which Assumption~\ref{ass:h} holds are all correct, so the whole primary curve can be read at once,
with no further correction for looking at many tolerances.

A simulation at the roster's scale, with \SimReplicates{} replicates per scenario under the real
design (Appendix~\ref{app:inference}), checks whether the guarantee is attained. Where
Assumption~\ref{ass:h} holds, no ordering was wrongly certified in any scenario, which bounds the
family-wise error, the chance that even one certified ordering is wrong, below \SimZeroFwerUpper{}
at \SimZeroFwerConfPct\% confidence; with heterogeneity
that the analysis does not allow for, the error rises to \SimFwerMisspecified{}, above the nominal
\AlphaAudit{}. When every model shares one underlying rate and the analysis uses $\varepsilon = 0$,
the two-verifier analysis errs at rate \SimFwerPropensityKTwo{} and the three-verifier analysis at
\SimFwerPropensityKThree{}, and a tolerance of \EpsPrereg{} removes the error in both. This is why we
do not lead with $\varepsilon = 0$.
